\textbf{Proof.}
Put \(\varepsilon=\varepsilon_{N_0,K_T,K_D}(\alpha)\), and let \(Z_r\), \(\mathfrak b_r\), and \(Q_r\) be the grid processes in the corrected linearization lemma. At every sufficiently large triple, that lemma gives
\[
 \sup_P P\{\|Z_r\|_{\infty,\mathcal G}>c_r^\xi+\omega_r/2\}
 \le\alpha-\varepsilon/4,
 \qquad
 \sup_P P\{\|Q_r\|_{\infty,\mathcal G}>\omega_r/2\}
 \le\varepsilon/4,
\]
and \(\|\mathfrak b_r\|_{\infty,\mathcal G}\le\mathbf 1\{r\ge2\}\eta_D(K_D)\). A union bound therefore gives, uniformly in \(P\),
\[
 P\!\left\{
 \max_{g\in\mathcal G}|\widehat\theta_r(g)-\theta(g)|
 \le c_r^\xi+\mathbf 1\{r\ge2\}\eta_D(K_D)+\omega_r
 \right\}\ge1-\alpha. \tag{1}
\]
This is a gridwise statement: no interpolation error has been included in \(Z_r\), \(\mathfrak b_r\), or \(Q_r\).

Set \(q=\min\{\beta+1,2\}\). The interpolation proposition gives
\[
 \|\widehat\theta_r-\theta\|_{\infty,\mathcal I}
 \le\max_{g\in\mathcal G}|\widehat\theta_r(g)-\theta(g)|
 +c_{\mathrm{int}}h^q,
 \qquad h^q\le\rho_r. \tag{2}
\]
Projection onto \([0,1]\) cannot increase distance from \(\theta(g)\in[0,1]\). Linear interpolation of the projected grid values has its maximum deviation from the linear interpolant of \(\theta\) at a grid endpoint, and \(\theta(0)=1\). The corrected radius in \(\mathrm{def:fpcr\mbox{-}multiplier\mbox{-}band}\) is exactly the right side of (1) plus the last term of (2). Hence (1)--(2) imply
\[
 P\{\widehat B_r^{\mathrm{lo}}(t)\le\theta(t)\le
 \widehat B_r^{\mathrm{up}}(t)\text{ for every }t\in\mathcal I\}
 \ge1-\alpha.
\]
This is stronger than \(1-\alpha-\varepsilon\), so \(\widehat B_r\in\mathcal H_r(N_0,K_T,K_D;\alpha)\).

Clipping endpoints to \([0,1]\) can only reduce width, and therefore
\[
 w(\widehat B_r)
 \le2c_r^\xi+2\mathbf 1\{r\ge2\}\eta_D(K_D)
 +2\omega_r+2c_{\mathrm{int}}h^q. \tag{3}
\]
For all sufficiently large triples, \(q_{N_0,K_T,K_D}\le1-\alpha/2\). The bounded martingale envelopes and the entropy calculation in the linearization proof give a uniform conditional maximal second moment of order \(\rho_r^2\). Markov's inequality at the fixed tail probability \(\alpha/2\) yields
\[
 \sup_{P\in\mathfrak P_r}E_Pc_r^\xi\le C\rho_r. \tag{4}
\]
Also \(h^q\le\rho_r\), and the channel-by-channel calculation in the linearization lemma gives \(\omega_r=o(\rho_r)\) uniformly on every three-dimensional tail. Taking expectations in (3), using (4), and increasing \(a_0\) if necessary yield
\[
 \sup_{P\in\mathfrak P_r}E_Pw(\widehat B_r)
 \le C\{\rho_r+\mathbf 1\{r\ge2\}\eta_D(K_D)\}
 =C\bar\rho_r. \tag{5}
\]
At \(r=1\), the risk-exhaustion term is absent by definition, so \(\bar\rho_1=\rho_1\). If \(r\ge2\) and \(\eta_D(K_D)=O(\rho_r)\), then \(\bar\rho_r=O(\rho_r)\). No step compares the relative growth of \(K_T\), \(N_0\), and \(K_D\). \(\square\)